
Or
prenons pour $X$ une \uline{variété de drapeaux}
\[
 X=G/B,
\]
où $G$ est un groupe algébrique semi-simple simplement connexe sur le corps algébriquement clos $k$, et $B$ un sous-groupe de Borel [5]. Lorsque $k=\C$, alors la décomposition cellulaire de Bruhat de $X$ [5, Exp. 13] implique que l'homomorphisme canonique
\[
 (*)\qquad K^\bullet(X)\longrightarrow K(X(\C))
\]
est un isomorphisme, de même d'ailleurs que l'homomorphisme canonique
\[
 (**)\qquad A(X)\longrightarrow H^*(X(\C),\Z),
\]
où $X(\C)$ est l'espace compact des points complexes de $X$, muni de sa topologie habituelle déduite de la topologie de $\C$. D'autre part, il a été prouvé par L. HODGKIN [12] que $K(X(\C))$ est bien engendré par les classes des faisceaux inversibles sur $X$. Il en est donc de même de $K^\bullet(X)$ lui-même, et il n'est d'ailleurs pas difficile d'en déduire, par une technique de spécialisation utilisant les groupes semi-simples déployés sur $\Z$ (SGA 3), que le même résultat reste valable pour tout corps $k$ algébriquement clos. D'ailleurs, la décomposition cellulaire de $X$ (ou la bijectivité de $(**)$ dans le cas $k=\C$) implique aisément que $A(X)$ est sans torsion, donc grâce à 4.2 l'homomorphisme $\phi:A(X)\longrightarrow\Gtop^\bullet(X)$ est un isomorphisme. De ces faits il résulte que $j$ est un isomorphisme si et seulement si $A(X)$ est engendré par ses éléments de degré $1$, ou encore, si $k=\C$, si et seulement si la cohomologie entière de $X(\C)$ est engendrée par les éléments de degré $2$, i.e. si le classifiant $B_{G(\C)}$ est sans torsion. Il est connu d'autre part [15, p. 217 th. B (1)$\Leftrightarrow$(2), et 2.5] que ceci est le cas si et seulement si $G$ est "sans torsion", ou encore si $G$ est isomorphe à un produit de groupes $\operatorname{S\ell}(n)$ et $\operatorname{Sp}(m)$. Il suffit donc de prendre pour $G$ un groupe $\operatorname{SO}(n)$ avec $n\geq 7$, ou un groupe simple exceptionnel, pour avoir un exemple où $j$ n'est pas un isomorphisme.

\textbf{4.6.} On notera que dans cet exemple, il est hors de doute que le "bon" anneau est l'anneau sans torsion $A(X)\simeq\Gtop^\bullet(X)$, ayant comme base les cellules de la décomposition de Bruhat, anneau également isomorphe à la cohomologie entière de $X(\C)$ si $k=\C$; par contre, l'anneau $\operatorname{Gr}^\bullet(X)$ apparaît comme un anneau ayant de la torsion de nature assez pathologique. Cet exemple confirme donc la conception suivant laquelle l'invariant $\operatorname{Gr}^\bullet(X)$ n'est guère
